\booksection{ภาคผนวก: การวิเคราะห์การขยายหน่วยความจำมือถือแอนดรอยด์}

\chapterintro{เวลาซื้อโทรศัพท์ คุณมักเห็นโฆษณาแบบ “12GB+8GB RAM” เหมือนได้ RAM เพิ่มมา 8GB ฟรีๆ แต่มันทำให้โทรศัพท์เร็วขึ้นจริงไหม? หรืออาจทำร้ายโทรศัพท์? ในส่วนนี้เราจะอธิบายเทคโนโลยี “การขยายหน่วยความจำ” ให้ชัดเจน}

\booksubsection{นิยามพื้นฐาน}

การ “ขยายหน่วยความจำ” ของโทรศัพท์แอนดรอยด์ (เรียกอย่างอื่นว่า ส่วนขยายหน่วยความจำ, การหลอมรวมหน่วยความจำ, RAM Plus ฯลฯ)\textbf{ไม่ได้เพิ่มชิป RAM จริงทางกายภาพ}แต่เป็นเทคนิคระดับซอฟต์แวร์ของระบบ: นำพื้นที่ว่างบางส่วนในตัวจัดเก็บของเครื่อง (ROM/UFS แฟลช) มาจัดสรรแบบไดนามิก จำลองให้เป็นหน่วยความจำเสมือน (RAM) ใช้งาน

ลองเปรียบเทียบ: RAM จริงคือ “โต๊ะทำงานความเร็วสูง” ส่วนหน่วยความจำขยายคือ “ห้องเก็บของชั่วคราว” เวลาโต๊ะทำงานไม่พอ ระบบจะพาข้อมูลที่ยังไม่ได้ใช้ไปเก็บไว้ที่ห้องเก็บของ เวลาต้องใช้อีกก็ย้ายกลับมา ห้องเก็บของกว้างขวางแต่ขนไปขนมาแน่นอนว่าช้ากว่าวางไว้บนโต๊ะทำงานโดยตรง

\booksubsection{หลักการทางเทคนิค}

วิธีที่ผู้ผลิตกระแสหลักนำมาใช้มีหลักๆ 2 แบบ:

\subsubsection*{บีบอัดแบบ zRAM}

แบ่งพื้นที่ส่วนหนึ่งภายใน RAM มาเป็นพื้นที่บีบอัดข้อมูล นำโค้ดของแอปที่ไม่ได้ใช้งานมาจัดเก็บแบบบีบอัด อัตราการบีบอัดราว 2:1 ข้อมูลยังอยู่ใน RAM ความเร็วจึงเร็ว ไม่ทำให้แฟลชเสื่อม แต่ต้องให้ CPU ช่วยบีบอัดและคลาย

\subsubsection*{พาร์ติชัน swap}

จัดพื้นที่คงที่ใน UFS แฟลชมาเป็นหน่วยความจำเสมือนโดยตรง เขียนข้อมูลของโปรเซสพื้นหลังลงไปเก็บ พึ่งพาแบนด์วิดท์ของตัวจัดเก็บ ความเร็วอ่านเขียนจึงต่ำกว่า RAM จริงมาก

ค่า “12GB+8GB” ที่ผู้ผลิตโฆษณาโดยทั่วไปคือการรวมกันของ RAM จริง 12GB + หน่วยความจำเสมือนบนแฟลช 8GB ที่น่าสังเกตคือ zRAM ตามค่าเริ่มต้นไม่กินพื้นที่จัดเก็บในเครื่อง จะไปสร้างไฟล์สลับบนตัวจัดเก็บก็ต่อเมื่อ zRAM ถูกข้อมูลเติมจนเต็มเท่านั้น ซึ่งกรณีแบบนี้เกิดขึ้นน้อยมาก

\booksubsection{ช่องว่างความเร็วเป็นคอขวดสำคัญ}

ข้อจำกัดโดยธรรมชาติของเทคนิคนี้คือช่องว่างของความเร็ว UFS 4.0 กระแสหลักปัจจุบันอ่านต่อเนื่องได้ราว 4200MB/s ซึ่งเป็นเพียง 1/10–1/15 ของ RAM จริงอย่าง LPDDR5X เท่านั้น ในช่วงโหลดเบาของ zRAM ความเร็วจึงใกล้เคียง RAM จริง แต่พอเข้าสู่สถานการณ์โหลดหนัก ความเร็วจะดิ่งลงทันทีเหลือเพียงหลายสิบต่อหนึ่ง

\booksubsection{สถานการณ์ที่เหมาะสมและขอบเขตประโยชน์}

หน่วยความจำขยายไม่ใช่รอบด้าน ผลลัพธ์ขึ้นกับรุ่นเครื่องและสถานการณ์ใช้งาน:

\begin{itemize}
  \item \textbf{กรณีที่มีประโยชน์}: รุ่นที่มี RAM จริง 6GB + UFS 3.1 ขึ้นไป พื้นที่ว่างเหลือ $\geq$30GB เมื่อเปิดขยาย 3GB จะรักษาแอปพื้นหลังไว้ได้อีก 3-4 ตัว อัตราโหลดหนักลดลง 68\% เหมาะกับผู้ใช้งานหนักหลายงานพร้อมกัน เช่น ไรเดอร์ส่งอาหาร คนขับรถรับส่งผ่านแอป
  \item \textbf{การใช้งานประจำวัน}: ต่อการเปิด WeChat หลายบัญชี การรักษาแท็บเว็บไว้ การตอบข้อความแชทเรียลไทม์ ดีขึ้นชัดเจน แอปพื้นหลังเปิดใหม่ตั้งแต่ต้นใช้เวลาสั้นลงชัดเจน
  \item \textbf{กรณีที่ได้ผลจำกัด}: สถานการณ์ที่ต้องการแบนด์วิดท์สูงอย่างการโหลดเกมใหญ่ การตัดต่อวิดีโอ 4K ได้รับการยกระดับแบบจำกัด
  \item \textbf{กรณีที่มีโทษ}: รุ่นที่มีแรม 8GB + UFS 2.2 และพื้นที่ว่างไม่พอ เปิดใช้แล้วเฟรมเกมแกว่งมากขึ้นกว่าเดิม ดูวิดีโอสั้นกินไฟมากขึ้น สลับแอปกระตุก
  \item \textbf{กรณีที่ไม่มีผล}: รุ่นแรม 16GB ขึ้นไปได้ประโยชน์เกือบเป็นศูนย์ รักษาแอปพื้นหลังเพิ่มได้มากที่สุด 1 ตัว แต่ทำให้ภาระ CPU เพิ่ม 12\% ร้อนขึ้น 2-3℃ สำหรับเครื่องที่ RAM เกิน 12GB แนะนำให้ปิดทิ้งเลย เพื่อไม่ต้องเสียการคำนวณบีบอัดโดยไร้ความหมาย
\end{itemize}

\booksubsection{การตั้งชื่อของแต่ละแบรนด์และวิธีเปิดใช้งาน}

แต่ละแบรนด์เรียกไม่เหมือนกัน: Xiaomi เรียก “ส่วนขยายหน่วยความจำ” Huawei เรียก “การหลอมรวมหน่วยความจำ” Samsung เรียก “RAM Plus” vivo ตั้งชื่อว่า “Extended RAM” OnePlus เรียกว่า “Memory Extension” ฟังก์ชันนี้เป็นความสามารถดั้งเดิมของระบบ ไม่ได้ผูกตรงๆ กับเวอร์ชันพื้นฐานของแอนดรอยด์

เส้นทางเปิดทั่วไปคือ: เข้า [ตั้งค่า] [การตั้งค่าเพิ่มเติม] หาหัวข้อย่อย [หน่วยความจำ] หรือ [ประสิทธิภาพ] เปิด “การขยายหน่วยความจำ” และเลือกระดับความจุ (มัก 2GB–7GB) รีสตาร์ทแล้วจึงมีผล เมื่อเปิดแล้วระบบจะแสดงหน่วยความจำรวมเป็นแบบ “8GB+4GB”

\begin{tcbwarn}{ข้อควรระวัง}
\begin{itemize}
  \item การขยายหน่วยความจำจะกินพื้นที่จัดเก็บในเครื่องบางส่วน โปรดมั่นใจว่าพื้นที่ว่างเหลือไม่น้อยกว่า 20-30GB มิฉะนั้นอาจกระทบความลื่นไหลของระบบ
  \item การใช้แฟลชเป็นหน่วยความจำเสมือนมีความเสื่อมจากการอ่านเขียน แม้ UFS ยุคปัจจุบันจะมีอายุการใช้งานยาว แต่การสลับข้อมูลถี่สูงต่อเนื่องนานๆ ก็ยังอาจทำให้ตัวจัดเก็บเสื่อมเร็วขึ้น
  \item ถ้าโทรศัพท์มี RAM จริงน้อยอยู่แล้ว (เช่น 4GB หรือต่ำกว่า) การเปิดขยายหน่วยความจำทำได้แค่ผ่อนความเดือดร้อนชั่วคราว ไม่สามารถแก้ปัญหากระตุกได้จากราก
\end{itemize}

\end{tcbwarn}

\begin{tcbinfo}{สรุป}
การขยายหน่วยความจำเป็นเทคนิคปรับแต่งระดับระบบที่โตมาแล้ว ช่วยยกระดับความสามารถหลายงานพร้อมกันและอัตรารักษาแอปพื้นหลังของเครื่องระดับกลาง-ล่างได้โดยฮาร์ดแวร์ไม่เปลี่ยน แต่มันถูกจำกัดด้วยหุบเขาความเร็วมโหฬารระหว่างแฟลชกับ RAM จริง ไม่ใช่ “ยาวิเศษรักษาได้ทุกโรค” จะเปิดหรือไม่ เปิดเท่าไร ต้องพิจารณาจากขนาด RAM ของรุ่นเครื่อง สเปกแฟลช พื้นที่ว่างที่เหลือ และนิสัยการใช้งานประกอบกัน

\end{tcbinfo}

\begin{tcbpractice}{ลงมือ: ตรวจดูว่าโทรศัพท์คุณมีการขยายหน่วยความจำหรือไม่}
\begin{enumerate}
  \item เปิด [ตั้งค่า] ค้นหา "ส่วนขยายหน่วยความจำ" "การหลอมรวมหน่วยความจำ" "RAM Plus" หรือ "Extended RAM"
  \item พบแล้วดูว่าตอนนี้เปิดใช้อยู่หรือไม่ และตั้งค่าความจุขยายไว้เท่าไร
  \item เทียบกับเนื้อหาในส่วนนี้แล้วประเมิน: RAM จริงของเครื่องคุณเท่าไร? สเปกแฟลชเป็นอะไร? พื้นที่จัดเก็บว่างพอหรือไม่?
  \item ลองปิดหรือลดความจุขยายลง แล้วสังเกตการรักษาแอปพื้นหลังและความร้อนตลอดหนึ่งสัปดาห์ เพื่อหาค่าตั้งที่เหมาะกับตัวเองที่สุด
\end{enumerate}

\end{tcbpractice}
